\chapter{Normative Systems: Logic, Probability and Rational Choice}
\label{sec:p11-normative-systems-comparison}

Chater and Oaksford's Oxford 2012 chapter compares three kinds of consistency:
relations among beliefs, among degrees of belief, and among preferences,
beliefs and choices. Its subject is neither a single universal score of
rationality nor a merger of three interchangeable calculi
\cite{ChaterOaksford2012NormativeSystems}. This treatment preserves that
comparison, repairs the stated model example, and makes the assumptions
behind the probabilistic and choice sketches explicit. The formal examples
and constructions below have full proofs; they are not participant data.

\section{Three roles for normative theory}
\label{sec:p11-nsc-roles}

The source separates the psychological question of how thought works from
the normative question of what counts as appropriate reasoning. Normative
theory can nevertheless contribute to psychology in three distinct ways.
First, it can specify a benchmark: asking whether an answer is deductively
correct already presupposes an account of deduction. Second, it can suggest
a descriptive mechanism, as arithmetic suggests a possible account of a
calculator's operations. This is controversial: a normative result can be
approximated by different mechanisms or not implemented at all. Third,
some philosophical accounts use rationality constraints to make belief,
intention and goal attribution possible. Random utterances or movements
do not identify those attitudes without an interpretive account
\cite[pp.~11--12]{ChaterOaksford2012NormativeSystems}.

The third role is the source's Davidson/Quine-inspired position, not an
empirical theorem that every departure makes a person meaningless. Its
closing food example shows the inference problem: saying that food is
safe while refusing to eat can complicate belief attribution. Refusal
could also reflect another preference, uncertainty or circumstance;
one behavior does not uniquely reveal a belief.

\begin{hbdefinition}{def:p11-nsc-evaluands}{Keep the evaluated objects distinct}
Logical consistency asks whether a represented set of propositions can
jointly be true; consequence asks what follows if its premises hold.
Probability coherence concerns a specified system of numerical credences
over events. Choice consistency concerns comparisons and selections on
a declared alternative domain. A normative benchmark, a hypothesized
mental process, and an interpretation of observed conduct are distinct
uses of a formalism.

These consistency standards do not, by themselves, fix true premises,
appropriate prior beliefs, valuable ends, or the normative authority of
conforming to a standard. They supply conditional constraints at their
specified scope, not every fact needed to decide how to think or act.
\end{hbdefinition}

Alien-abduction beliefs and willing a sports team to win are the source's
background-belief illustrations. Those claims can fit a stipulated view
of alien activity or psychic powers while conflicting with other commitments.
Their physicalist/psychic-power contrast concerns commitments that actually
exclude one another under the intended interpretation, not a definition
declaring every unusual claim inconsistent. The formal consistency question
is not the question whether that background
is evidentially justified. Likewise, a choice to smoke or exercise cannot
be assessed from its label alone without beliefs, consequences and values.
The examples do not endorse their background assumptions or provide
health advice \cite[p.~12]{ChaterOaksford2012NormativeSystems}.

\section{The logical example: consequence and satisfiability}
\label{sec:p11-nsc-logic}

\begin{hbconvention}{def:p11-nsc-logic-model}{The source's local first-order notation}
Use a nonempty domain, unary predicates \(\mathrm{worm}\) and
\(\mathrm{warble}\), and constant \(\mathrm{Albert}\), whose denotation
is \(a\). Their extensions are \(S_{\rm worm}\) and
\(S_{\rm warble}\). The source writes \(\&\) for conjunction and
\(\subset\) for inclusion allowing equality; those local meanings are
retained here. The conditional \(\to\) is material, not automatically
every ordinary-language conditional.

Let \(\Gamma\) consist of
\(\forall x(\mathrm{worm}(x)\to\mathrm{warble}(x))\) and
\(\mathrm{worm}(\mathrm{Albert})\), and write
\(\varphi=\mathrm{warble}(\mathrm{Albert})\).
The source's three-belief set is \(\Gamma\cup\{\neg\varphi\}\).
\end{hbconvention}

Universal instantiation yields that being a worm implies warbling for
Albert; conditional elimination yields warbling from Albert's being a
worm. Together with its negation, conjunction introduction makes the
contradiction explicit. These are applications of the classical calculus,
not new editorial inference-rule displays
\cite[pp.~12--14]{ChaterOaksford2012NormativeSystems}.

\begin{hbtheorem}{thm:p11-nsc-logical-checks}{proved=annex:p11.nsc.logical-checks}{The example, its repairs and the consistency test}
The two premises \(\Gamma\) derive and semantically entail \(\varphi\).
The three-belief set has no model. Deleting any one of its three members
produces a satisfiable set; consequently inconsistency alone does not
select which member to discard.

For the corrected positive example, use domain \(\{3,4,5,8\}\),
\(a=5\), \(S_{\rm worm}=\{3,5,8\}\), and
\(S_{\rm warble}=\{3,4,5,8\}\). These satisfy the universal premise,
the worm premise and positive warbling, not negative warbling.

More generally, for classical sentence semantics,
\(\Gamma\entails\varphi\) exactly when
\(\Gamma\cup\{\neg\varphi\}\) has no model. An unsatisfiable
premise set semantically entails every sentence. For the registered
classical natural-deduction calculus, the corresponding refutation test
uses negation and classical reductio under their actual discharge conditions.
\end{hbtheorem}

The source's comparison of proof theory and model theory distinguishes
deriving a contradiction from the absence of a satisfying interpretation.
Soundness excludes deriving a contradiction from satisfiable premises;
completeness supplies a contradiction when premises are unsatisfiable.
The equivalent consequence formulations compare derivability with truth
in every premise model. The actual first-order soundness and strong
completeness interfaces are \ref{thm:fol-soundness} and
\ref{thm:fol-completeness}. The latter assumes Choice for a general
signature. The countable-signature corollary
\ref{cor:fol-completeness-countable} makes no use of Choice and covers
the finite-signature example above.

Semantic completeness of a calculus is different from completeness of
a chosen arithmetic axiom theory, and from completeness for truth in
one intended structure. Expressing arithmetic in a first-order language
does not invalidate first-order semantic completeness. The source's
incompleteness discussion must preserve that distinction. Arithmetic-theory
incompleteness and truth in the intended natural-number structure are
different questions from the first-order semantic adequacy considered here.

\section{Static constraints are not a complete revision policy}
\label{sec:p11-nsc-revision}

If the first two worm beliefs are protected, accepting their consequence
removes the conflict with doubt about warbling. If the doubt is protected,
one might instead reconsider the generalization or Albert's classification.
The models above demonstrate these different coherent repairs. Classical
logic does not assign their evidential or practical priorities. The source
thus denies that logic alone is a complete theory of how to change beliefs
\cite[p.~15]{ChaterOaksford2012NormativeSystems}.

The locality claim needs its interpretation fixed. Truth preservation in
a represented argument does not depend on the zoological meaning of its
predicate names. But deciding how language should be represented or which
premises to accept can require background knowledge. Understanding a news
story or conversation commonly involves defeasible additions absent from
the spoken premises; treating those additions as certain changes the task.

The source contrasts mental-logics approaches with mental-model approaches,
and reviews logical-development and symbolic-AI accounts. These are
theoretical/historical positions, not a new experiment or evidence that
correct responses identify a unique architecture. Deontic, modal, temporal
and second-order logics have different languages and semantics. Their
fusion is not guaranteed by putting their names together. Ordinary
conditionals, in particular, admit competing formal interpretations
\cite[pp.~13--15]{ChaterOaksford2012NormativeSystems}.

\section{Probability coherence and the evidence inversion}
\label{sec:p11-nsc-probability}

\begin{hbdefinition}{def:p11-nsc-probability-model}{A declared finite credence model}
Let \(\Omega\) be finite and nonempty. A probability measure \(\Pr\)
assigns nonnegative masses totaling one, with event probabilities obtained
by summing their masses. Thus disjoint-event probabilities add,
\(\Pr(\Omega)=1\), and complements have probability one minus the
original probability. This formal coherence is distinct from calibration
to observed outcomes or correctness of the modeled alternatives.

For events \(A,B\), the source's joint notation \(\Pr(A,B)\) means
\(\Pr(A\cap B)\). If \(\Pr(B)>0\), define
\(\Pr(A\mid B)=\Pr(A,B)/\Pr(B)\). A null conditioning event
does not acquire a conditional probability by dividing zero by zero.
\end{hbdefinition}

\begin{hbtheorem}{thm:p11-nsc-bayes}{proved=annex:p11.nsc.bayes}{The source's Bayes identity with its domain}
If \(\Pr(A)>0\) and \(\Pr(B)>0\), then
\[
 \Pr(B\mid A)=\frac{\Pr(A\mid B)\Pr(B)}{\Pr(A)}.
\]
Knowing \(\Pr(A\mid B)\) alone does not fix the converse. Coherent
models can share that likelihood and have different posterior values.
The formula is an identity within the declared law, not by itself a
universal command to revise every belief after uncertain testimony.
\end{hbtheorem}

The chapter uses a subjective interpretation: probabilities model degrees
of belief. It distinguishes that interpretation and Bayesian modeling
from the elementary validity of Bayes' formula, which does not depend
on endorsing Bayesian psychology. Its references to Cox, Kolmogorov and
Ramsey concern different foundations; the chapter gives no single common
hypothesis set proving a uniquely obligatory calculus from all meanings
of belief. Graphical models and their computational properties are
identified as important but explicitly beyond its treatment
\cite[p.~16]{ChaterOaksford2012NormativeSystems}.

If the information learned is that a positive-probability event occurred,
conditioning a retained model is a specified update. If the evidence is
only a fallible report, one conditions on the report event, not silently
on the reported proposition. A changing model or uncertain new evidence
requires further assumptions. The full finite conditional law and
hypothesis-partition formulas are in \ref{thm:p15-finite-conditioning-laws}
and \ref{thm:p15-finite-bayes}. Neither coherent probabilities nor a good
fit to responses proves that people performed that calculation mentally.

\section{The choice conditions and formal utility}
\label{sec:p11-nsc-choice}

The source's alternatives \(P,M,D\) are now a poke in the ribs, a
restaurant meal and a decoy. They are alternatives, not propositions or
probabilities. Preferring \(P\) to \(M\) in isolation does not contradict
an ordering merely because an observer dislikes that preference.
The issue is how selections across menus relate.

\begin{hbdefinition}{def:p11-nsc-menu-model}{A finite specification of contraction and expansion}
Let \(X\) be a finite nonempty option set. For every nonempty menu
\(B\subseteq X\), let \(C(B)\) be a nonempty subset of \(B\),
retaining all selected ties. \emph{Contraction} requires that, if
\(A\subseteq B\) and \(x\in A\cap C(B)\), then \(x\in C(A)\).
\emph{Expansion} requires that \(x\in C(A)\cap C(B)\) imply
\(x\in C(A\cup B)\). These conditions are quantified over nonempty
menus in this finite formalization of the source's examples.

Define revealed weak comparison by
\(x\succeq_C y\) exactly when \(x\in C(\{x,y\})\).
Transitivity of this relation is an additional condition. An ordinal
utility \(u\) represents it when weak comparison is exactly comparison
of its real values; utility here is a formal representation, not a
measure of moral worth or objective usefulness.
\end{hbdefinition}

\begin{hbtheorem}{thm:p11-nsc-menu-representation}{proved=annex:p11.nsc.menu-representation}{Finite menu consistency and representation}
Under contraction and expansion, \(\succeq_C\) is reflexive and complete,
and each \(C(B)\) is exactly the set of options weakly at least as good
as every member of \(B\). If revealed comparison is also transitive,
there is a real ordinal representation \(u\) and
\(C(B)=\operatorname{argmax}_{x\in B}u(x)\).
Conversely the maximal correspondence of a finite complete transitive
relation has contraction and expansion and recovers that relation.
\end{hbtheorem}

Choosing \(P\) from \(\{P,M,D\}\) but \(M\) from \(\{P,M\}\)
violates contraction. Choosing \(P\) from both \(\{P,M\}\) and
\(\{P,D\}\), but not from their union, violates expansion. The
finite theorem gives the source's sketch precise premises; arbitrary
ordered domains need further conditions for real representation.
The full weak-order, tied-option and ordinal-scale interfaces are in
Chapter~\ref{sec:p17-finite-preferences}. An increasing change of ordinal
scale preserves rankings but need not preserve averages of prize values
\cite[pp.~16--17]{ChaterOaksford2012NormativeSystems}.

\section{The money-pump argument is conditional}
\label{sec:p11-nsc-pump}

The source argues that an adversary could exploit the contraction pattern
by offering a trade from \(P\) to \(M\) in the pair menu, then from
\(M\) to \(P\) in the expanded menu, charging small fees. Its
additional practical premises matter: a free menu choice does not by
itself establish willingness to pay, unchanged preferences after losing
wealth, available trades or an agent's inability to recognize the cycle.

\begin{hbtheorem}{thm:p11-nsc-pump}{proved=annex:p11.nsc.pump}{Two-step loss and a refusal counterpolicy}
If both source trades are accepted, with positive fees
\(\Delta_1,\Delta_2\), the agent returns to \(P\) with wealth reduced
by \(\Delta_1+\Delta_2\). If the same affordable trades are accepted
for \(k\) rounds, the loss is \(k(\Delta_1+\Delta_2)\).

Contraction violation alone does not imply a positive loss. A policy
that makes the source's free-menu selections but refuses every transaction
reducing wealth exhibits the violation and pays no positive fee.
\end{hbtheorem}

This retains the normative warning without claiming that every menu
dependence drains every person's money. The counterpolicy need not be
endorsed as a complete rationality ideal; it establishes the missing
implication in the proposed argument. A preference change caused by
learning or different outcomes would also require a different model
\cite[p.~17]{ChaterOaksford2012NormativeSystems}.

\section{Lotteries require more than ordinal prize ranking}
\label{sec:p11-nsc-lotteries}

The source's substitution principle compares lotteries with the same
remaining outcome: replacing a component by a preferred one should
improve the lottery when its weight \(q\) is positive. At \(q=0\)
the lotteries coincide, so strict improvement is impossible. Its continuity
example seeks a weight \(s\) making an intermediate prize \(R\)
indifferent to a mixture of a better \(P\) and worse \(M\).
An ordinal ordering of prizes alone does not specify such weights.

For a precise finite-prize statement, use a complete transitive weak
preference on the \emph{whole lottery simplex}, mixture independence,
and continuity, as defined in Chapter~\ref{sec:p17-finite-expected-utility}.
The full representation theorem \ref{thm:p17-finite-vnm} then supplies
a consequence utility whose probability-weighted sum represents preference.
Its proof covers calibration, component replacement and all-indifferent
cases. The scale is unique only up to positive affine change in the
nondegenerate case, as proved in \ref{thm:p17-vnm-affine-uniqueness}.
These are the base/domain conditions absent from an unqualified reading
of the chapter's two-additional-conditions sketch
\cite[p.~17]{ChaterOaksford2012NormativeSystems}.

Expected-utility representation does not choose the values for a person,
prove a moral duty to maximize them, or establish that humans implement
the calculation. The chapter describes EU as both a benchmark and a
starting point for descriptive models, and distinguishes uncertain
outcomes from strategic interaction. Game theory is acknowledged but
not developed in its individual-choice treatment. Its historical reviews
of psychology, economics and behavioral ecology supply no new common
participant sample, measured efficacy or universal mechanism.

\section{One case, three different questions}
\label{sec:p11-nsc-worked}

\begin{hbdefinition}{def:p11-nsc-worked-model}{A fallible report about Albert}
The logical conclusion \(\varphi\) still follows \emph{if} both
premises \(\Gamma\) are true. Separately let \(H\) be the event that
Albert warbles and \(E\) a fallible positive report of the premises.
Declare a four-outcome probability law with joint masses for
\((H,E),(H,E^c),(H^c,E),(H^c,E^c)\) of
\(9/100,1/100,9/100,81/100\), respectively.
Reports can be wrong; their occurrence is not the truth of \(\Gamma\).

For a further declared decision, a correct classification has utility 1
and an incorrect one utility 0. Available acts are guessing that \(H\)
holds, guessing that it does not, or using a perfect test whose utility
cost is \(c\geq0\). Test accuracy and this utility model are stipulated;
they are not science about worms or an observed human decision.
\end{hbdefinition}

\begin{hbtheorem}{thm:p11-nsc-worked}{proved=annex:p11.nsc.worked}{Entailment, posterior weight and a cost-sensitive decision}
In the declared model, \(\Pr(H)=1/10\),
\(\Pr(E\mid H)=9/10\), \(\Pr(E\mid H^c)=1/10\), and
\(\Pr(H\mid E)=1/2\). After the report, each guess has expected
utility \(1/2\), while the perfect test has expected utility \(1-c\).
Testing is uniquely best when \(c<1/2\); the guesses are both best
when \(c>1/2\); at equality all three acts tie. Thus cost .2 favors
testing and cost .6 favors guessing at the same posterior probability.
\end{hbtheorem}

There is no contradiction between the conditional logical consequence
and the posterior half: the probability model does not assert certainty
of the premises. The same probability model supports different decisions
under different costs. Conversely, a favorable outcome cannot repair
an invalid consequence claim. This constructed case answers the lookup
question without merging truth preservation, credence and action value.

\section{World knowledge, inconsistency and computational limits}
\label{sec:p11-nsc-limits}

The source's first challenge is world knowledge. Understanding a simple
story can require a network of physical and social assumptions. Categories
such as chair, country, person and belief are not automatically defined
by a few axioms. The frame problem asks what changes and what remains
after an action, illustrated by an animator deciding what to redraw.
The authors' fractal/web metaphors identify a representational challenge;
they are not proofs that all such formalization is impossible.

The second challenge is inconsistency: choosing a repair and preventing
uncontrolled inference are different problems. Classical semantic explosion
was proved above; a paraconsistent alternative requires its own specified
consequence relation. Probability and choice failures should not be called
the same formal explosion without defining the relevant representation.
The source's low-crash-risk/fear-of-flying and saving/compulsive-spending
illustrations do not alone establish contradictory propositions: emotions,
losses, goals, commitments and constraints must be represented. Its Fermat
example illustrates the difficulty of identifying consequences, not a
promised Fermat proof or a newly measured inconsistency rate
\cite[pp.~18--19]{ChaterOaksford2012NormativeSystems}.

\begin{hbdefinition}{def:p11-nsc-sat-task}{The finite consistency benchmark}
\(\mathsf{SAT}\) asks whether a finitely encoded Boolean formula has
a satisfying valuation. Use the verifier characterization of
\(\mathsf{NP}\): an input \(x\), of length \(n\), belongs to a
language if some polynomially bounded certificate \(w\) is accepted
by a fixed deterministic polynomial-time verifier \(V\).
A polynomial witness-length bound is \(p(n)\). In the annex construction,
\(T\) is a polynomial time bound, \(L\) the initial tape length,
\(b\) a fixed cell-code width, \(\mathsf F\) the finite local transition
function, and \(z,a_g,b_g\) gate-output/input Boolean variables.
A polynomial-time many-one reduction constructs one output instance
whose answer is yes exactly when the original answer is yes.
\end{hbdefinition}

\begin{hbtheorem}{thm:p11-nsc-sat-complete}{proved=annex:p11.nsc.sat-complete}{The source's SAT benchmark, with a complete construction}
\(\mathsf{SAT}\) is \(\mathsf{NP}\)-complete under polynomial-time
many-one reductions. The construction reduces every verifier language
to satisfiability of a conjunction of clauses of at most three literals,
so this restricted clause problem is already \(\mathsf{NP}\)-complete.
\end{hbtheorem}

NP-completeness is not a proof that P differs from NP, that every instance
is hard, or that each human consistency task is infeasible. Unsatisfiability
and entailment have complementary decision formulations; they must not
be mislabeled simply because SAT is NP-complete. For unrestricted encoded
first-order validity, the full c.e.-completeness interface is
\ref{thm:p04-dh-validity-complete}; by negation, unrestricted satisfiability
cannot have a universal decider either. This does not apply indiscriminately
to every fragment or bounded example. Cost claims about probability and
choice also need the actual input representation and requested computation
\cite[pp.~19--20]{ChaterOaksford2012NormativeSystems}.

The source's closing transitive-preference and disjunctive-syllogism
examples likewise require the other comparisons or premises to remain
fixed: an accepted disjunction with a denied second disjunct entails the
first, but static logic does not prohibit revising the disjunction instead.
The source contrasts rejecting normative models as literal mechanisms
with using cheap approximations to them. A finite expected-value calculation
does not require infinite resources merely because an unrestricted task
can be difficult. Neither a complexity theorem nor a good response fit
identifies how a person reasons. Similarity, analogy, creativity and insight
are also not exhausted by static coherence, although the chapter leaves
open whether further normative principles could illuminate them.

\section{Selection and the neighboring philosophical interfaces}
\label{sec:p11-nsc-use}

\begin{hbmethod}{def:p11-nsc-selection}{Select the actual normative question}
Specify the representation, evaluated object and conclusion. For deduction,
record which premises are held fixed and which semantics/calculus applies.
For probability, specify events, the prior law, the evidence and the
conditioning domain. For choice, specify menus or lotteries, comparisons,
values, costs and available acts. State whether the goal is a benchmark,
a descriptive mechanism hypothesis, an interpretation of attitudes, or
a rule for revision. Check both success and failure cases under those
premises. Do not turn a fitted response, a coherent ranking, or a valid
conditional conclusion into evidence for the other two standards.
\end{hbmethod}

The full P2, P15 and P17 treatments give the source's mathematical
interfaces rather than a replacement source presentation. The actual
Knauff--Spohn introduction in Chapter~\ref{sec:p11-rationality-intro}
additionally distinguishes capacities, performances and broader standards.
Broome's \ref{def:p11-reasons-normative-senses} separates a correctness
standard from having a normative reason to conform; his
\ref{def:p11-reasons-roles} separates motivating from normative reasons.
The present consistency account does not settle those questions by calling
a utility scale rational. Its source is a philosophical/theoretical
comparison, not a training experiment or a universal justification of
every end, prior, inference system or mental architecture.
